\section{Sensores e sistemas de percepção}

Os sensores oferecem precisão e operam em condições adversas.
Os mais comuns no contexto espacial são o LiDAR (Light Detection and Ranging), câmeras estéreo, câmeras de alta resolução, sistemas GNSS e pseudolites.
Por fim, a fusão desses sensores possibilita um modelo situacional robusto e confiável.

\subsection{LiDAR (Light Detection and Ranging)}

De acordo com \cite{inpeLidar}, o LiDAR é um sensor remoto ativo, atuando de forma direta para captura de dados.
Este sensor possui sua própria fonte de energia (fonte de luz), o laser.
O LiDAR emite feixes de laser na banda do infravermelho próximo (IV) e é capaz de modelar a superfície do terreno de forma tridimensional., possibilitando gerar produtos como o Modelo Digital de Terreno e o Modelo Digital de Superfície que representam o terreno (sem nenhuma cobertura) e a superfície (edifícios, árvores etc), respectivamente.
A técnica LiDAR é utilizada principalmente para levantamentos topográficos, caracterização da estrutura da vegetação, bem como a volumetria de edificações e ambientes urbanos de forma mais rápida e confiável.

\subsubsection{O LiDAR em aplicações espaciais}

De acordo com \cite{winker1996}, o experimento \emph{Lidar In-space Technology Experiment (LITE)} foi desenvolvido pelo \emph{NASA Langley Research Center} e lançado a bordo do ônibus espacial \emph{Discovery} em setembro de 1994, como parte da missão STS-64.
Seu objetivo era validar tecnologias de LiDAR para aplicações espaciais, investigando sua capacidade de detectar camadas de aerossóis e nuvens na atmosfera terrestre.
Diferentemente dos sensores passivos tradicionalmente usados para observações atmosféricas, o LiDAR espacial oferece medições mais detalhadas e diretas de estrutura de nuvens e aerossóis, mesmo em condições noturnas ou sobre terrenos irregulares.

A missão coletou mais de 40 GB de dados em alta resolução da estrutura vertical da atmosfera.
Também validou a tecnologia; o experimento destacou os benefícios da utilização do LiDAR em futuras plataformas orbitais autônomas, ampliando sua utilização para estudos climáticos, previsão meteorológica e monitoramento ambiental.
O desenvolvimento do \emph{LITE} foi precedido por décadas de pesquisa em sistemas LiDAR aerotransportados e terrestres, que demonstraram sua eficácia na medição de vapor d'água, aerossóis e ozônio.
A maturidade da tecnologia permitiu que em 1985 o projeto fosse iniciado com o objetivo de validar a operação de um sistema LiDAR no espaço.
Em 1988, foi formado um grupo de coordenação científica para definir os requisitos do instrumento e planejar a missão.
O LITE demonstrou que sensores ativos baseados em laser podem complementar sistemas passivos e radares para estudos atmosféricos, pavimentando o caminho para novas gerações de satélites equipados com essa tecnologia \cite{winker1996}.
